% part: intuitionistic-logic
% chapter: introduction
% section: axiomatic-derivations

\documentclass[../../../include/open-logic-section]{subfiles}

\begin{document}

\olfileid{int}{int}{axd}

\olsection{\usetoken{P}{derivation} اصل‌موضوعی}

!!{derivation}های اصل‌موضوعی برای منطق گزاره‌ای شهودگرایانه، از نظر
مفهومی ساده‌ترین و از نظر تاریخی نخستین دستگاه‌های !!{derivation}
هستند. آن‌ها درست همانند منطق گزاره‌ای کلاسیک عمل می‌کنند.

\begin{defn}[!!^{derivability}]
اگر~$\Gamma$ مجموعه‌ای از !!{formula}های~$\Lang L$ باشد، آنگاه یک
\emph{!!{derivation}} از~$\Gamma$ دنبالهٔ متناهی
$!A_1$، \dots،~$!A_n$ از !!{formula}هاست که برای هر~$i \le n$ یکی از
شرایط زیر برقرار است:
\begin{enumerate}
\item $!A_i \in \Gamma$؛ یا
\item $!A_i$ یک اصل موضوع است؛ یا
\item $!A_i$ با وضع مقدم از دو~$!A_j$ و~$!A_k$، که~$j < i$ و~$k < i$،
  نتیجه می‌شود؛ یعنی~$!A_k \ident !A_j \lif !A_i$.
\end{enumerate}
\end{defn}

\begin{defn}[اصول موضوع]
مجموعهٔ~$\PAx$ از \emph{اصول موضوع} منطق گزاره‌ای شهودگرایانه شامل همهٔ
!!{formula}های صورت‌های زیر است:
\begin{align}
  & (!A \land !B) \lif !A \ollabel{ax:land1}\\
  & (!A \land !B) \lif !B \ollabel{ax:land2}\\
  & !A \lif (!B \lif (!A \land !B)) \ollabel{ax:land3}\\
  & !A \lif (!A \lor !B) \ollabel{ax:lor1}\\
  & !A \lif (!B \lor !A) \ollabel{ax:lor2}\\
  & (!A \lif !C) \lif ((!B \lif !C) \lif ((!A \lor !B) \lif !C)) \ollabel{ax:lor3}\\
  & !A \lif (!B \lif !A) \ollabel{ax:lif1}\\
  & (!A \lif (!B \lif !C)) \lif ((!A \lif !B) \lif (!A \lif !C)) \ollabel{ax:lif2}\\
  & \lfalse \lif !A \ollabel{ax:lfalse1}
\end{align}
\end{defn}

\begin{defn}[!!^{derivability}]
!!^a{formula}~$!A$ از~$\Gamma$ \emph{!!{derivable}} است، که می‌نویسیم
$\Gamma \Proves !A$، اگر !!a{derivation}ای از~$\Gamma$ وجود داشته
باشد که به~$!A$ پایان یابد.
\end{defn}

\begin{defn}[قضایا]
!!^a{formula}~$!A$ یک \emph{قضیه} است هرگاه !!a{derivation}ای از~$!A$
از مجموعهٔ تهی وجود داشته باشد. اگر~$!A$ قضیه باشد می‌نویسیم
$\Proves !A$، و اگر نباشد~$\Proves/ !A$.
\end{defn}

\begin{prop}
  اگر~$\Gamma \Proves !A$ در دستگاه !!{derivation} اصل‌موضوعی برای
  منطق شهودگرایانه، آنگاه~$\Gamma \Proves !A$ در منطق کلاسیک نیز.
  به‌ویژه، اگر~$!A$ قضیه‌ای شهودگرایانه باشد، قضیه‌ای کلاسیک نیز هست.
\end{prop}

\begin{proof}
  هر اصل موضوع شهودگرایانه یک اصل موضوع کلاسیک نیز هست؛ پس هر
  !!{derivation} در منطق شهودگرایانه !!a{derivation}ای در منطق کلاسیک
  نیز هست.
\end{proof}

\end{document}
